\documentclass[11pt,a4paper]{article}
\usepackage[margin=2.2cm]{geometry}
\usepackage[utf8]{inputenc}
\usepackage[T1]{fontenc}
\usepackage{amsmath,amssymb,bm}
\usepackage{xcolor}
\usepackage{booktabs}
\usepackage{titlesec}
\usepackage{enumitem}

\setlength{\parindent}{0pt}
\setlength{\parskip}{0.6em}

\titleformat{\section}{\Large\bfseries}{\thesection.}{0.5em}{}
\titleformat{\subsection}{\large\bfseries}{\thesubsection}{0.5em}{}

\title{\vspace{-1.5cm}\bfseries How the GRPO Training Loss Is Defined,\\
How It's Actually Computed, and Why the Numerical Value Sits Near Zero}
\author{}
\date{}

\begin{document}
\maketitle
\vspace{-1.5em}

\section{Loss definition}

For a single trajectory $i$ inside group $g$ (sampled with the current policy
$\pi_\theta$), the policy-gradient loss is
\begin{equation}
\boxed{\;\;
L_{g,i} \;=\; -\;A_{g,i}\;\cdot\;\overline{\log p_\theta}\!\left(\,\text{actions of trajectory }i\,\right)
\;\;}
\end{equation}

\textbf{First factor --- the advantage} $A_{g,i}$ --- is the within-group
$z$-score of the terminal reward:
\begin{equation}
A_{g,i} \;=\; \frac{r_{g,i} \;-\; \bar r_g}{\sigma_{r,g} \;+\; \varepsilon}
\qquad
\text{where}\;\; \bar r_g = \tfrac{1}{G}\sum_{j} r_{g,j},\;\;
\sigma_{r,g} = \sqrt{\tfrac{1}{G}\sum_j (r_{g,j} - \bar r_g)^2}.
\end{equation}

In words: ``this trajectory's terminal reward, minus the average reward of
the eight trajectories in its group, divided by the within-group reward
standard deviation.''

\textbf{Second factor --- the mean log-probability} $\overline{\log p_\theta}$
--- is the policy's log-probability of the trajectory's actually-taken
action tokens, averaged over those tokens:
\begin{equation}
\overline{\log p_\theta} \;=\; \frac{1}{T}\sum_{t=1}^{T} \log p_\theta\!\left(\text{token}_t \;\big|\; \text{prefix}_t\right).
\end{equation}

In words: ``during the rollout you actually emitted tokens
$\text{token}_1, \dots, \text{token}_T$. Now run the model on the
re-tokenized chat and ask: what log-probability does the (current) policy
give to each of those exact tokens? Take the average.''

\textbf{Total loss} is the average of $L_{g,i}$ over all trajectory items
in the batch:
\begin{equation}
L_{\text{total}} \;=\; \frac{1}{n_{\text{items}}}\,\sum_{g,\,i} L_{g,i}.
\end{equation}

\section{How it is actually computed (in code)}

See \texttt{scripts/17\_grpo\_static.py}. The compute path for one item is:

\begin{enumerate}[leftmargin=2em]
\item Concatenate \texttt{prefix} (system+user history) and the trajectory's
\texttt{assistant\_text}, run a forward pass through the model:
\begin{verbatim}
out    = model(input_ids=enc["input_ids"], attention_mask=...)
logits = out.logits[:, :-1, :]    # shape (1, T-1, V), V = 152K vocab
\end{verbatim}

\item For every token position, compute $-\log p_\theta(\text{label}_t)$ via
chunked cross-entropy (chunked along the time axis to avoid materializing
the full $(T \times V)$ log-softmax tensor):
\begin{verbatim}
nll = F.cross_entropy(logits, labels, reduction="none")
# nll[t] = - log p_theta(token_t | prefix_t)
\end{verbatim}

\item Mask out the prefix region, average over the assistant region:
\begin{verbatim}
mean_logp = -(nll * assistant_mask).sum() / assistant_mask.sum()
\end{verbatim}

\item Form the per-item PG loss and \emph{immediately} call
\texttt{.backward()} (must not accumulate
$L_{\text{total}} \mathrel{+}= L_b$ across items --- that holds every
forward pass's autograd graph in memory and OOMs):
\begin{verbatim}
loss_i = -(advantage * mean_logp) / n_items
loss_i.backward()
\end{verbatim}
After the loop completes, call \texttt{optimizer.step()} once.
\end{enumerate}

\section{Why the loss numerically sits near zero}
\label{sec:near_zero}

The cleanest way to see this is with concrete numbers. Take a single group
of $G = 8$ trajectories with terminal rewards
\[
\mathbf{r} \;=\; [\,1,\,1,\,0,\,0,\,1,\,0,\,1,\,0\,].
\]
Then
\begin{align}
\bar r &= 0.5,\\
\sigma_r &= 0.5,\\
A_i &= \frac{r_i - 0.5}{0.5} \;=\; [\,+1,\,+1,\,-1,\,-1,\,+1,\,-1,\,+1,\,-1\,].
\end{align}

\textbf{Crucial observation: $\;\sum_i A_i \;=\; 0\,$.}\;\;
This is forced by the $z$-score normalization (the numerator subtracts the
mean, so the standardized values sum to zero by construction).

For a freshly initialized 7B policy, the per-token log-probability of
short ``Thought:..\textbackslash nAction:..'' continuations is around $-3$
(a vocabulary of 152K and slightly biased token frequencies). Suppose the
eight trajectories happen to have slightly different mean log-probs:
\[
\overline{\log p_\theta} \;\approx\; [\,-3.1,\,-2.9,\,-3.0,\,-3.2,\,-2.8,\,-3.1,\,-2.9,\,-3.0\,].
\]

Then $\text{loss}_i = -A_i \cdot \overline{\log p_\theta}_i$:
\begin{center}
\begin{tabular}{cccr}
\toprule
$i$ & $A_i$ & $\overline{\log p_\theta}_i$ & $\text{loss}_i$ \\
\midrule
1 & $+1$ & $-3.1$ & $+3.1$ \\
2 & $+1$ & $-2.9$ & $+2.9$ \\
3 & $-1$ & $-3.0$ & $-3.0$ \\
4 & $-1$ & $-3.2$ & $-3.2$ \\
5 & $+1$ & $-2.8$ & $+2.8$ \\
6 & $-1$ & $-3.1$ & $-3.1$ \\
7 & $+1$ & $-2.9$ & $+2.9$ \\
8 & $-1$ & $-3.0$ & $-3.0$ \\
\midrule
\multicolumn{3}{r}{$\sum_i$} & $-0.6$ \\
\multicolumn{3}{r}{$\div\,n_{\text{items}}\!=\!64$} & $\boxed{\approx -0.009}$ \\
\bottomrule
\end{tabular}
\end{center}

This number is exactly in the range we observed empirically
($[-0.01,\,+0.06]$ across 20 training steps).

\section{The intuition}

Each individual term $\text{loss}_i$ has magnitude $\sim 3$ (because
$|A_i| \approx 1$ and $|\overline{\log p_\theta}_i| \approx 3$).
\textbf{But $\sum_i A_i = 0$ kills the bulk of those $\pm 3$
contributions almost completely.}

What survives is the correlation between $A_i$ and the \emph{deviation} of
$\overline{\log p_\theta}_i$ from its group mean:
\begin{equation}
L_{\text{total}} \;\approx\; \sum_i A_i \cdot \Bigl(\,\overline{\log p_\theta}_i \;-\; \mathrm{mean}\bigl(\overline{\log p_\theta}\bigr)\,\Bigr).
\label{eq:residual}
\end{equation}

That correlation is the genuine training signal --- and it is small in
magnitude precisely because the policy at any given moment doesn't yet
distinguish good from bad trajectories all that strongly. Hence
$L_{\text{total}} \sim 0.01$ in absolute value. \textbf{This is not a bug;
it is the definitional behaviour of GRPO.}

\section{Why the loss is not expected to monotonically decrease}

\begin{itemize}[leftmargin=1.6em]
\item \textbf{Supervised-learning loss} measures ``how wrong is the
prediction'' --- it is bounded below by zero and decreases monotonically as
the model learns.

\item \textbf{Policy-gradient loss} (REINFORCE / GRPO) is constructed so
that its \emph{gradient} is an unbiased estimator of $\nabla J(\theta)$,
the expected reward gradient. Its \emph{value} is just an artefact: by the
group-centering $\sum_i A_i = 0$, the loss has expected value $0$. It
oscillates around $0$ regardless of whether the policy is improving.
\end{itemize}

A useful analogy: a supervised loss is like ``how far you are from the
finish line'' --- decreasing means you are getting closer. A
policy-gradient loss is like ``which way is the wind blowing'' --- the
direction carries information, but the wind \emph{speed} does not have
to fall to zero as you sail.

\section{What you should look at instead of the loss}

To check whether GRPO training is doing useful work, the right diagnostics
are:

\begin{enumerate}[leftmargin=1.6em]
\item \textbf{LoRA gradient $\bm{L^2}$ norm trend}: should slowly decay as
the policy becomes more confident in its preferred actions. Across our
20 steps it bounces in $[0.05,\,0.20]$ without a clear trend --- meaning
``the trainer is producing real gradients'' but ``too few steps to see
convergence dynamics.''

\item \textbf{Held-out task reward / success rate}: roll out the
\emph{updated} model on a held-out task suite and measure win rate.
\textbf{We have not done this yet}, so any claim that the model is
``getting better at ALFWorld'' is currently unsupported.
\end{enumerate}

\section{What the experiment as it stands does and does not show}

\begin{center}
\begin{tabular}{lc}
\toprule
\textbf{Claim} & \textbf{Evidence} \\
\midrule
The loss is computed correctly (forward, backward, optimizer all run) & \checkmark \\
The gradient signal magnitude is preserved (baseline 0.086, gated 0.098) & \checkmark \\
The gate saves 32.2\% wall-clock per training step & \checkmark \\
The model becomes better at ALFWorld & \emph{not measured} \\
The training run is converged & \emph{no} (20 steps; loss is also not the right indicator) \\
\bottomrule
\end{tabular}
\end{center}

\end{document}
